\QHeading{第1题　口算。}
\par\noindent


\begin{oralgrid}
\OralItem{1}{3.7+0.6}{4.3}
\OralItem{2}{2.52+1.47}{3.99}
\OralItem{3}{7.28+2.72}{10}
\OralItem{4}{2.7-1.9}{0.8}
\OralItem{5}{3.69-2.78}{0.91}
\OralItem{6}{7.95-3.45}{4.5}
\OralItem{7}{3.6+2.5}{6.1}
\OralItem{8}{5.42+6.27}{11.69}\end{oralgrid}
\par

\QHeading{口算思路补充（对应上表题号）}

\Answer{上面列出得数，下面按小题给出不同计算思路。计数单位法适用于小数，整数题可比较凑整、拆分或商不变等方法。}
\MethodTitle{第1小题：$3.7+0.6$}

\MethodBlock{口算思路一：按相同计数单位计算}{先看成若干个十分之一，按整数加减，最后还原成小数。}{\begin{aligned}[t]&\phantom{={}}3.7+0.6\\
&=(37+6)\div 10\\
&=43\div 10\\
&=4.3\end{aligned}}
\MethodTitle{第2小题：$2.52+1.47$}

\MethodBlock{口算思路一：按相同计数单位计算}{先看成若干个百分之一，按整数加减，最后还原成小数。}{\begin{aligned}[t]&\phantom{={}}2.52+1.47\\
&=(252+147)\div 100\\
&=399\div 100\\
&=3.99\end{aligned}}
\MethodTitle{第3小题：$7.28+2.72$}

\MethodBlock{口算思路一：按相同计数单位计算}{先看成若干个百分之一，按整数加减，最后还原成小数。}{\begin{aligned}[t]&\phantom{={}}7.28+2.72\\
&=(728+272)\div 100\\
&=1000\div 100\\
&=10\end{aligned}}

\MethodBlock{口算思路二：凑整再减去多加的数}{把第二个数凑成整数，再减去多加的部分。}{\begin{aligned}[t]&\phantom{={}}7.28+2.72\\
&=7.28+3-0.28\\
&=10.28-0.28\\
&=10\end{aligned}}
\MethodTitle{第4小题：$2.7-1.9$}

\MethodBlock{口算思路一：按相同计数单位计算}{先看成若干个十分之一，按整数加减，最后还原成小数。}{\begin{aligned}[t]&\phantom{={}}2.7-1.9\\
&=(27-19)\div 10\\
&=8\div 10\\
&=0.8\end{aligned}}

\MethodBlock{口算思路二：补差法}{从减数先补到一个整数，再补到被减数，两段差相加。}{\begin{aligned}[t]&\phantom{={}}2.7-1.9\\
&=(2-1.9)+(2.7-2)\\
&=0.1+0.7\\
&=0.8\end{aligned}}
\MethodTitle{第5小题：$3.69-2.78$}

\MethodBlock{口算思路一：按相同计数单位计算}{先看成若干个百分之一，按整数加减，最后还原成小数。}{\begin{aligned}[t]&\phantom{={}}3.69-2.78\\
&=(369-278)\div 100\\
&=91\div 100\\
&=0.91\end{aligned}}

\MethodBlock{口算思路二：补差法}{从减数先补到一个整数，再补到被减数，两段差相加。}{\begin{aligned}[t]&\phantom{={}}3.69-2.78\\
&=(3-2.78)+(3.69-3)\\
&=0.22+0.69\\
&=0.91\end{aligned}}
\MethodTitle{第6小题：$7.95-3.45$}

\MethodBlock{口算思路一：按相同计数单位计算}{先看成若干个百分之一，按整数加减，最后还原成小数。}{\begin{aligned}[t]&\phantom{={}}7.95-3.45\\
&=(795-345)\div 100\\
&=450\div 100\\
&=4.5\end{aligned}}

\MethodBlock{口算思路二：补差法}{从减数先补到一个整数，再补到被减数，两段差相加。}{\begin{aligned}[t]&\phantom{={}}7.95-3.45\\
&=(4-3.45)+(7.95-4)\\
&=0.55+3.95\\
&=4.5\end{aligned}}
\MethodTitle{第7小题：$3.6+2.5$}

\MethodBlock{口算思路一：按相同计数单位计算}{先看成若干个十分之一，按整数加减，最后还原成小数。}{\begin{aligned}[t]&\phantom{={}}3.6+2.5\\
&=(36+25)\div 10\\
&=61\div 10\\
&=6.1\end{aligned}}
\MethodTitle{第8小题：$5.42+6.27$}

\MethodBlock{口算思路一：按相同计数单位计算}{先看成若干个百分之一，按整数加减，最后还原成小数。}{\begin{aligned}[t]&\phantom{={}}5.42+6.27\\
&=(542+627)\div 100\\
&=1169\div 100\\
&=11.69\end{aligned}}


\par\noindent

\QHeading{第2题　列竖式计算，最后一行要验算。}

\Problem{（1）$50-14.14=35.86$}
\par


\MethodTitle{方法一：按题意列竖式}
\Answer{\begin{tikzpicture}[x=1em,y=1em,baseline=(current bounding box.center),draw=answerred,text=answerred]
\node[anchor=base,inner sep=0pt] at (0.000,0) {$5$};
\node[anchor=base,inner sep=0pt] at (0.800,0) {$0$};
\node[anchor=base,inner sep=0pt] at (1.600,0) {$.$};
\node[anchor=base,inner sep=0pt] at (1.900,0) {$0$};
\node[anchor=base,inner sep=0pt] at (2.700,0) {$0$};
\node[anchor=base,inner sep=0pt] at (0.000,-1.1) {$1$};
\node[anchor=base,inner sep=0pt] at (0.800,-1.1) {$4$};
\node[anchor=base,inner sep=0pt] at (1.600,-1.1) {$.$};
\node[anchor=base,inner sep=0pt] at (1.900,-1.1) {$1$};
\node[anchor=base,inner sep=0pt] at (2.700,-1.1) {$4$};
\node[anchor=base,inner sep=0pt] at (0.000,-2.3) {$3$};
\node[anchor=base,inner sep=0pt] at (0.800,-2.3) {$5$};
\node[anchor=base,inner sep=0pt] at (1.600,-2.3) {$.$};
\node[anchor=base,inner sep=0pt] at (1.900,-2.3) {$8$};
\node[anchor=base,inner sep=0pt] at (2.700,-2.3) {$6$};
\node[anchor=base,inner sep=0pt] at (-.7,-1.1) {$-$};
\draw[line width=.5pt] (-1,-1.52)--(3.300,-1.52);
\draw[line width=.35pt] (-0.160,.03)--(0.160,.43);
\node[anchor=base,inner sep=0pt,font=\normalsize] at (0.000,.72) {$4$};
\draw[line width=.35pt] (0.640,.03)--(0.960,.43);
\node[anchor=base,inner sep=0pt,font=\normalsize] at (0.800,.72) {$9$};
\draw[line width=.35pt] (1.740,.03)--(2.060,.43);
\node[anchor=base,inner sep=0pt,font=\normalsize] at (1.900,.72) {$9$};
\draw[line width=.35pt] (2.540,.03)--(2.860,.43);
\node[anchor=base,inner sep=0pt,font=\normalsize] at (2.700,.72) {$10$};
\end{tikzpicture}}

\MethodBlock{方法二：拆成整数部分和小数部分}{上方竖式及题目要求的验算仍保留；本法用于核对。}{\begin{aligned}[t]&\phantom{={}}50-14.14\\
&=50-14-0.14\\
&=36-0.14\\
&=35.86\end{aligned}}

\MethodBlock{方法三：按相同计数单位计算}{上方竖式及题目要求的验算仍保留；本法用于核对。先看成若干个百分之一，按整数加减，最后还原成小数。}{\begin{aligned}[t]&\phantom{={}}50-14.14\\
&=(5000-1414)\div 100\\
&=3586\div 100\\
&=35.86\end{aligned}}
\Problem{（2）$53.69+14.86=68.55$}

\MethodTitle{方法一：按题意列竖式}
\Answer{\begin{tikzpicture}[x=1em,y=1em,baseline=(current bounding box.center),draw=answerred,text=answerred]
\node[anchor=base,inner sep=0pt] at (0.000,0) {$5$};
\node[anchor=base,inner sep=0pt] at (0.800,0) {$3$};
\node[anchor=base,inner sep=0pt] at (1.600,0) {$.$};
\node[anchor=base,inner sep=0pt] at (1.900,0) {$6$};
\node[anchor=base,inner sep=0pt] at (2.700,0) {$9$};
\node[anchor=base,inner sep=0pt] at (0.000,-1.1) {$1$};
\node[anchor=base,inner sep=0pt] at (0.800,-1.1) {$4$};
\node[anchor=base,inner sep=0pt] at (1.600,-1.1) {$.$};
\node[anchor=base,inner sep=0pt] at (1.900,-1.1) {$8$};
\node[anchor=base,inner sep=0pt] at (2.700,-1.1) {$6$};
\node[anchor=base,inner sep=0pt] at (0.000,-2.3) {$6$};
\node[anchor=base,inner sep=0pt] at (0.800,-2.3) {$8$};
\node[anchor=base,inner sep=0pt] at (1.600,-2.3) {$.$};
\node[anchor=base,inner sep=0pt] at (1.900,-2.3) {$5$};
\node[anchor=base,inner sep=0pt] at (2.700,-2.3) {$5$};
\node[anchor=base,inner sep=0pt] at (-.7,-1.1) {$+$};
\draw[line width=.5pt] (-1,-1.52)--(3.300,-1.52);
\node[anchor=base,inner sep=0pt,font=\normalsize] at (1.900,.72) {$1$};
\node[anchor=base,inner sep=0pt,font=\normalsize] at (0.800,.72) {$1$};
\end{tikzpicture}}

\MethodBlock{方法二：凑整再减去多加的数}{上方竖式及题目要求的验算仍保留；本法用于核对。把第二个数凑成整数，再减去多加的部分。}{\begin{aligned}[t]&\phantom{={}}53.69+14.86\\
&=53.69+15-0.14\\
&=68.69-0.14\\
&=68.55\end{aligned}}

\MethodBlock{方法三：拆成整数部分和小数部分}{上方竖式及题目要求的验算仍保留；本法用于核对。}{\begin{aligned}[t]&\phantom{={}}53.69+14.86\\
&=53.69+14+0.86\\
&=67.69+0.86\\
&=68.55\end{aligned}}
\Problem{（3）$53.62+12.59=66.21$}

\MethodTitle{方法一：按题意列竖式}
\Answer{\begin{tikzpicture}[x=1em,y=1em,baseline=(current bounding box.center),draw=answerred,text=answerred]
\node[anchor=base,inner sep=0pt] at (0.000,0) {$5$};
\node[anchor=base,inner sep=0pt] at (0.800,0) {$3$};
\node[anchor=base,inner sep=0pt] at (1.600,0) {$.$};
\node[anchor=base,inner sep=0pt] at (1.900,0) {$6$};
\node[anchor=base,inner sep=0pt] at (2.700,0) {$2$};
\node[anchor=base,inner sep=0pt] at (0.000,-1.1) {$1$};
\node[anchor=base,inner sep=0pt] at (0.800,-1.1) {$2$};
\node[anchor=base,inner sep=0pt] at (1.600,-1.1) {$.$};
\node[anchor=base,inner sep=0pt] at (1.900,-1.1) {$5$};
\node[anchor=base,inner sep=0pt] at (2.700,-1.1) {$9$};
\node[anchor=base,inner sep=0pt] at (0.000,-2.3) {$6$};
\node[anchor=base,inner sep=0pt] at (0.800,-2.3) {$6$};
\node[anchor=base,inner sep=0pt] at (1.600,-2.3) {$.$};
\node[anchor=base,inner sep=0pt] at (1.900,-2.3) {$2$};
\node[anchor=base,inner sep=0pt] at (2.700,-2.3) {$1$};
\node[anchor=base,inner sep=0pt] at (-.7,-1.1) {$+$};
\draw[line width=.5pt] (-1,-1.52)--(3.300,-1.52);
\node[anchor=base,inner sep=0pt,font=\normalsize] at (1.900,.72) {$1$};
\node[anchor=base,inner sep=0pt,font=\normalsize] at (0.800,.72) {$1$};
\end{tikzpicture}}

\MethodBlock{方法二：拆成整数部分和小数部分}{上方竖式及题目要求的验算仍保留；本法用于核对。}{\begin{aligned}[t]&\phantom{={}}53.62+12.59\\
&=53.62+12+0.59\\
&=65.62+0.59\\
&=66.21\end{aligned}}

\MethodBlock{方法三：按相同计数单位计算}{上方竖式及题目要求的验算仍保留；本法用于核对。先看成若干个百分之一，按整数加减，最后还原成小数。}{\begin{aligned}[t]&\phantom{={}}53.62+12.59\\
&=(5362+1259)\div 100\\
&=6621\div 100\\
&=66.21\end{aligned}}
\Problem{（4）$12.98-3.8=9.18$}

\MethodTitle{方法一：按题意列竖式}
\Answer{\begin{tikzpicture}[x=1em,y=1em,baseline=(current bounding box.center),draw=answerred,text=answerred]
\node[anchor=base,inner sep=0pt] at (0.000,0) {$1$};
\node[anchor=base,inner sep=0pt] at (0.800,0) {$2$};
\node[anchor=base,inner sep=0pt] at (1.600,0) {$.$};
\node[anchor=base,inner sep=0pt] at (1.900,0) {$9$};
\node[anchor=base,inner sep=0pt] at (2.700,0) {$8$};
\node[anchor=base,inner sep=0pt] at (0.800,-1.1) {$3$};
\node[anchor=base,inner sep=0pt] at (1.600,-1.1) {$.$};
\node[anchor=base,inner sep=0pt] at (1.900,-1.1) {$8$};
\node[anchor=base,inner sep=0pt] at (2.700,-1.1) {$0$};
\node[anchor=base,inner sep=0pt] at (0.800,-2.3) {$9$};
\node[anchor=base,inner sep=0pt] at (1.600,-2.3) {$.$};
\node[anchor=base,inner sep=0pt] at (1.900,-2.3) {$1$};
\node[anchor=base,inner sep=0pt] at (2.700,-2.3) {$8$};
\node[anchor=base,inner sep=0pt] at (-.7,-1.1) {$-$};
\draw[line width=.5pt] (-1,-1.52)--(3.300,-1.52);
\draw[line width=.35pt] (-0.160,.03)--(0.160,.43);
\node[anchor=base,inner sep=0pt,font=\normalsize] at (0.000,.72) {$0$};
\draw[line width=.35pt] (0.640,.03)--(0.960,.43);
\node[anchor=base,inner sep=0pt,font=\normalsize] at (0.800,.72) {$12$};
\end{tikzpicture}}

\MethodBlock{方法二：多减再加回}{上方竖式及题目要求的验算仍保留；本法用于核对。减去凑成的整数后，把多减的部分加回来。}{\begin{aligned}[t]&\phantom{={}}12.98-3.8\\
&=12.98-4+0.2\\
&=8.98+0.2\\
&=9.18\end{aligned}}

\MethodBlock{方法三：拆成整数部分和小数部分}{上方竖式及题目要求的验算仍保留；本法用于核对。}{\begin{aligned}[t]&\phantom{={}}12.98-3.8\\
&=12.98-3-0.8\\
&=9.98-0.8\\
&=9.18\end{aligned}}
\Problem{（5）$12.96-6.2=6.76$}

\MethodTitle{方法一：按题意列竖式}
\Answer{\begin{tikzpicture}[x=1em,y=1em,baseline=(current bounding box.center),draw=answerred,text=answerred]
\node[anchor=base,inner sep=0pt] at (0.000,0) {$1$};
\node[anchor=base,inner sep=0pt] at (0.800,0) {$2$};
\node[anchor=base,inner sep=0pt] at (1.600,0) {$.$};
\node[anchor=base,inner sep=0pt] at (1.900,0) {$9$};
\node[anchor=base,inner sep=0pt] at (2.700,0) {$6$};
\node[anchor=base,inner sep=0pt] at (0.800,-1.1) {$6$};
\node[anchor=base,inner sep=0pt] at (1.600,-1.1) {$.$};
\node[anchor=base,inner sep=0pt] at (1.900,-1.1) {$2$};
\node[anchor=base,inner sep=0pt] at (2.700,-1.1) {$0$};
\node[anchor=base,inner sep=0pt] at (0.800,-2.3) {$6$};
\node[anchor=base,inner sep=0pt] at (1.600,-2.3) {$.$};
\node[anchor=base,inner sep=0pt] at (1.900,-2.3) {$7$};
\node[anchor=base,inner sep=0pt] at (2.700,-2.3) {$6$};
\node[anchor=base,inner sep=0pt] at (-.7,-1.1) {$-$};
\draw[line width=.5pt] (-1,-1.52)--(3.300,-1.52);
\draw[line width=.35pt] (-0.160,.03)--(0.160,.43);
\node[anchor=base,inner sep=0pt,font=\normalsize] at (0.000,.72) {$0$};
\draw[line width=.35pt] (0.640,.03)--(0.960,.43);
\node[anchor=base,inner sep=0pt,font=\normalsize] at (0.800,.72) {$12$};
\end{tikzpicture}}

\MethodBlock{方法二：拆成整数部分和小数部分}{上方竖式及题目要求的验算仍保留；本法用于核对。}{\begin{aligned}[t]&\phantom{={}}12.96-6.2\\
&=12.96-6-0.2\\
&=6.96-0.2\\
&=6.76\end{aligned}}

\MethodBlock{方法三：按相同计数单位计算}{上方竖式及题目要求的验算仍保留；本法用于核对。先看成若干个百分之一，按整数加减，最后还原成小数。}{\begin{aligned}[t]&\phantom{={}}12.96-6.2\\
&=(1296-620)\div 100\\
&=676\div 100\\
&=6.76\end{aligned}}
\Problem{（6）$4.05+9.09=13.14$}

\MethodTitle{方法一：按题意列竖式}
\Answer{\begin{tikzpicture}[x=1em,y=1em,baseline=(current bounding box.center),draw=answerred,text=answerred]
\node[anchor=base,inner sep=0pt] at (0.800,0) {$4$};
\node[anchor=base,inner sep=0pt] at (1.600,0) {$.$};
\node[anchor=base,inner sep=0pt] at (1.900,0) {$0$};
\node[anchor=base,inner sep=0pt] at (2.700,0) {$5$};
\node[anchor=base,inner sep=0pt] at (0.800,-1.1) {$9$};
\node[anchor=base,inner sep=0pt] at (1.600,-1.1) {$.$};
\node[anchor=base,inner sep=0pt] at (1.900,-1.1) {$0$};
\node[anchor=base,inner sep=0pt] at (2.700,-1.1) {$9$};
\node[anchor=base,inner sep=0pt] at (0.000,-2.3) {$1$};
\node[anchor=base,inner sep=0pt] at (0.800,-2.3) {$3$};
\node[anchor=base,inner sep=0pt] at (1.600,-2.3) {$.$};
\node[anchor=base,inner sep=0pt] at (1.900,-2.3) {$1$};
\node[anchor=base,inner sep=0pt] at (2.700,-2.3) {$4$};
\node[anchor=base,inner sep=0pt] at (-.7,-1.1) {$+$};
\draw[line width=.5pt] (-1,-1.52)--(3.300,-1.52);
\node[anchor=base,inner sep=0pt,font=\normalsize] at (1.900,.72) {$1$};
\node[anchor=base,inner sep=0pt,font=\normalsize] at (0.000,.72) {$1$};
\end{tikzpicture}}

\MethodBlock{方法二：拆成整数部分和小数部分}{上方竖式及题目要求的验算仍保留；本法用于核对。}{\begin{aligned}[t]&\phantom{={}}4.05+9.09\\
&=4.05+9+0.09\\
&=13.05+0.09\\
&=13.14\end{aligned}}

\MethodBlock{方法三：按相同计数单位计算}{上方竖式及题目要求的验算仍保留；本法用于核对。先看成若干个百分之一，按整数加减，最后还原成小数。}{\begin{aligned}[t]&\phantom{={}}4.05+9.09\\
&=(405+909)\div 100\\
&=1314\div 100\\
&=13.14\end{aligned}}
\Problem{（7）$2.56+4.44=7$}

\MethodTitle{方法一：按题意列竖式}
\Answer{\begin{tikzpicture}[x=1em,y=1em,baseline=(current bounding box.center),draw=answerred,text=answerred]
\node[anchor=base,inner sep=0pt] at (0.000,0) {$2$};
\node[anchor=base,inner sep=0pt] at (0.800,0) {$.$};
\node[anchor=base,inner sep=0pt] at (1.100,0) {$5$};
\node[anchor=base,inner sep=0pt] at (1.900,0) {$6$};
\node[anchor=base,inner sep=0pt] at (0.000,-1.1) {$4$};
\node[anchor=base,inner sep=0pt] at (0.800,-1.1) {$.$};
\node[anchor=base,inner sep=0pt] at (1.100,-1.1) {$4$};
\node[anchor=base,inner sep=0pt] at (1.900,-1.1) {$4$};
\node[anchor=base,inner sep=0pt] at (0.000,-2.3) {$7$};
\node[anchor=base,inner sep=0pt] at (0.800,-2.3) {$.$};
\node[anchor=base,inner sep=0pt] at (1.100,-2.3) {$0$};
\node[anchor=base,inner sep=0pt] at (1.900,-2.3) {$0$};
\node[anchor=base,inner sep=0pt] at (-.7,-1.1) {$+$};
\draw[line width=.5pt] (-1,-1.52)--(2.500,-1.52);
\node[anchor=base,inner sep=0pt,font=\normalsize] at (1.100,.72) {$1$};
\node[anchor=base,inner sep=0pt,font=\normalsize] at (0.000,.72) {$1$};
\end{tikzpicture}\hspace{3pt}验算：\begin{tikzpicture}[x=1em,y=1em,baseline=(current bounding box.center),draw=answerred,text=answerred]
\node[anchor=base,inner sep=0pt] at (0.000,0) {$7$};
\node[anchor=base,inner sep=0pt] at (0.800,0) {$.$};
\node[anchor=base,inner sep=0pt] at (1.100,0) {$0$};
\node[anchor=base,inner sep=0pt] at (1.900,0) {$0$};
\node[anchor=base,inner sep=0pt] at (0.000,-1.1) {$4$};
\node[anchor=base,inner sep=0pt] at (0.800,-1.1) {$.$};
\node[anchor=base,inner sep=0pt] at (1.100,-1.1) {$4$};
\node[anchor=base,inner sep=0pt] at (1.900,-1.1) {$4$};
\node[anchor=base,inner sep=0pt] at (0.000,-2.3) {$2$};
\node[anchor=base,inner sep=0pt] at (0.800,-2.3) {$.$};
\node[anchor=base,inner sep=0pt] at (1.100,-2.3) {$5$};
\node[anchor=base,inner sep=0pt] at (1.900,-2.3) {$6$};
\node[anchor=base,inner sep=0pt] at (-.7,-1.1) {$-$};
\draw[line width=.5pt] (-1,-1.52)--(2.500,-1.52);
\draw[line width=.35pt] (-0.160,.03)--(0.160,.43);
\node[anchor=base,inner sep=0pt,font=\normalsize] at (0.000,.72) {$6$};
\draw[line width=.35pt] (0.940,.03)--(1.260,.43);
\node[anchor=base,inner sep=0pt,font=\normalsize] at (1.100,.72) {$9$};
\draw[line width=.35pt] (1.740,.03)--(2.060,.43);
\node[anchor=base,inner sep=0pt,font=\normalsize] at (1.900,.72) {$10$};
\end{tikzpicture}}

\MethodBlock{方法二：拆成整数部分和小数部分}{上方竖式及题目要求的验算仍保留；本法用于核对。}{\begin{aligned}[t]&\phantom{={}}2.56+4.44\\
&=2.56+4+0.44\\
&=6.56+0.44\\
&=7\end{aligned}}

\MethodBlock{方法三：按相同计数单位计算}{上方竖式及题目要求的验算仍保留；本法用于核对。先看成若干个百分之一，按整数加减，最后还原成小数。}{\begin{aligned}[t]&\phantom{={}}2.56+4.44\\
&=(256+444)\div 100\\
&=700\div 100\\
&=7\end{aligned}}
\Problem{（8）$28.3-15.06=13.24$}

\MethodTitle{方法一：按题意列竖式}
\Answer{\begin{tikzpicture}[x=1em,y=1em,baseline=(current bounding box.center),draw=answerred,text=answerred]
\node[anchor=base,inner sep=0pt] at (0.000,0) {$2$};
\node[anchor=base,inner sep=0pt] at (0.800,0) {$8$};
\node[anchor=base,inner sep=0pt] at (1.600,0) {$.$};
\node[anchor=base,inner sep=0pt] at (1.900,0) {$3$};
\node[anchor=base,inner sep=0pt] at (2.700,0) {$0$};
\node[anchor=base,inner sep=0pt] at (0.000,-1.1) {$1$};
\node[anchor=base,inner sep=0pt] at (0.800,-1.1) {$5$};
\node[anchor=base,inner sep=0pt] at (1.600,-1.1) {$.$};
\node[anchor=base,inner sep=0pt] at (1.900,-1.1) {$0$};
\node[anchor=base,inner sep=0pt] at (2.700,-1.1) {$6$};
\node[anchor=base,inner sep=0pt] at (0.000,-2.3) {$1$};
\node[anchor=base,inner sep=0pt] at (0.800,-2.3) {$3$};
\node[anchor=base,inner sep=0pt] at (1.600,-2.3) {$.$};
\node[anchor=base,inner sep=0pt] at (1.900,-2.3) {$2$};
\node[anchor=base,inner sep=0pt] at (2.700,-2.3) {$4$};
\node[anchor=base,inner sep=0pt] at (-.7,-1.1) {$-$};
\draw[line width=.5pt] (-1,-1.52)--(3.300,-1.52);
\draw[line width=.35pt] (1.740,.03)--(2.060,.43);
\node[anchor=base,inner sep=0pt,font=\normalsize] at (1.900,.72) {$2$};
\draw[line width=.35pt] (2.540,.03)--(2.860,.43);
\node[anchor=base,inner sep=0pt,font=\normalsize] at (2.700,.72) {$10$};
\end{tikzpicture}\hspace{3pt}验算：\begin{tikzpicture}[x=1em,y=1em,baseline=(current bounding box.center),draw=answerred,text=answerred]
\node[anchor=base,inner sep=0pt] at (0.000,0) {$1$};
\node[anchor=base,inner sep=0pt] at (0.800,0) {$3$};
\node[anchor=base,inner sep=0pt] at (1.600,0) {$.$};
\node[anchor=base,inner sep=0pt] at (1.900,0) {$2$};
\node[anchor=base,inner sep=0pt] at (2.700,0) {$4$};
\node[anchor=base,inner sep=0pt] at (0.000,-1.1) {$1$};
\node[anchor=base,inner sep=0pt] at (0.800,-1.1) {$5$};
\node[anchor=base,inner sep=0pt] at (1.600,-1.1) {$.$};
\node[anchor=base,inner sep=0pt] at (1.900,-1.1) {$0$};
\node[anchor=base,inner sep=0pt] at (2.700,-1.1) {$6$};
\node[anchor=base,inner sep=0pt] at (0.000,-2.3) {$2$};
\node[anchor=base,inner sep=0pt] at (0.800,-2.3) {$8$};
\node[anchor=base,inner sep=0pt] at (1.600,-2.3) {$.$};
\node[anchor=base,inner sep=0pt] at (1.900,-2.3) {$3$};
\node[anchor=base,inner sep=0pt] at (2.700,-2.3) {$0$};
\node[anchor=base,inner sep=0pt] at (-.7,-1.1) {$+$};
\draw[line width=.5pt] (-1,-1.52)--(3.300,-1.52);
\node[anchor=base,inner sep=0pt,font=\normalsize] at (1.900,.72) {$1$};
\end{tikzpicture}}

\MethodBlock{方法二：拆成整数部分和小数部分}{上方竖式及题目要求的验算仍保留；本法用于核对。}{\begin{aligned}[t]&\phantom{={}}28.3-15.06\\
&=28.3-15-0.06\\
&=13.3-0.06\\
&=13.24\end{aligned}}

\MethodBlock{方法三：按相同计数单位计算}{上方竖式及题目要求的验算仍保留；本法用于核对。先看成若干个百分之一，按整数加减，最后还原成小数。}{\begin{aligned}[t]&\phantom{={}}28.3-15.06\\
&=(2830-1506)\div 100\\
&=1324\div 100\\
&=13.24\end{aligned}}
\Problem{（9）$36.78+9.5=46.28$}

\MethodTitle{方法一：按题意列竖式}
\Answer{\begin{tikzpicture}[x=1em,y=1em,baseline=(current bounding box.center),draw=answerred,text=answerred]
\node[anchor=base,inner sep=0pt] at (0.000,0) {$3$};
\node[anchor=base,inner sep=0pt] at (0.800,0) {$6$};
\node[anchor=base,inner sep=0pt] at (1.600,0) {$.$};
\node[anchor=base,inner sep=0pt] at (1.900,0) {$7$};
\node[anchor=base,inner sep=0pt] at (2.700,0) {$8$};
\node[anchor=base,inner sep=0pt] at (0.800,-1.1) {$9$};
\node[anchor=base,inner sep=0pt] at (1.600,-1.1) {$.$};
\node[anchor=base,inner sep=0pt] at (1.900,-1.1) {$5$};
\node[anchor=base,inner sep=0pt] at (2.700,-1.1) {$0$};
\node[anchor=base,inner sep=0pt] at (0.000,-2.3) {$4$};
\node[anchor=base,inner sep=0pt] at (0.800,-2.3) {$6$};
\node[anchor=base,inner sep=0pt] at (1.600,-2.3) {$.$};
\node[anchor=base,inner sep=0pt] at (1.900,-2.3) {$2$};
\node[anchor=base,inner sep=0pt] at (2.700,-2.3) {$8$};
\node[anchor=base,inner sep=0pt] at (-.7,-1.1) {$+$};
\draw[line width=.5pt] (-1,-1.52)--(3.300,-1.52);
\node[anchor=base,inner sep=0pt,font=\normalsize] at (0.800,.72) {$1$};
\node[anchor=base,inner sep=0pt,font=\normalsize] at (0.000,.72) {$1$};
\end{tikzpicture}\hspace{3pt}验算：\begin{tikzpicture}[x=1em,y=1em,baseline=(current bounding box.center),draw=answerred,text=answerred]
\node[anchor=base,inner sep=0pt] at (0.000,0) {$4$};
\node[anchor=base,inner sep=0pt] at (0.800,0) {$6$};
\node[anchor=base,inner sep=0pt] at (1.600,0) {$.$};
\node[anchor=base,inner sep=0pt] at (1.900,0) {$2$};
\node[anchor=base,inner sep=0pt] at (2.700,0) {$8$};
\node[anchor=base,inner sep=0pt] at (0.800,-1.1) {$9$};
\node[anchor=base,inner sep=0pt] at (1.600,-1.1) {$.$};
\node[anchor=base,inner sep=0pt] at (1.900,-1.1) {$5$};
\node[anchor=base,inner sep=0pt] at (2.700,-1.1) {$0$};
\node[anchor=base,inner sep=0pt] at (0.000,-2.3) {$3$};
\node[anchor=base,inner sep=0pt] at (0.800,-2.3) {$6$};
\node[anchor=base,inner sep=0pt] at (1.600,-2.3) {$.$};
\node[anchor=base,inner sep=0pt] at (1.900,-2.3) {$7$};
\node[anchor=base,inner sep=0pt] at (2.700,-2.3) {$8$};
\node[anchor=base,inner sep=0pt] at (-.7,-1.1) {$-$};
\draw[line width=.5pt] (-1,-1.52)--(3.300,-1.52);
\draw[line width=.35pt] (-0.160,.03)--(0.160,.43);
\node[anchor=base,inner sep=0pt,font=\normalsize] at (0.000,.72) {$3$};
\draw[line width=.35pt] (0.640,.03)--(0.960,.43);
\node[anchor=base,inner sep=0pt,font=\normalsize] at (0.800,.72) {$15$};
\draw[line width=.35pt] (1.740,.03)--(2.060,.43);
\node[anchor=base,inner sep=0pt,font=\normalsize] at (1.900,.72) {$12$};
\end{tikzpicture}}

\MethodBlock{方法二：拆成整数部分和小数部分}{上方竖式及题目要求的验算仍保留；本法用于核对。}{\begin{aligned}[t]&\phantom{={}}36.78+9.5\\
&=36.78+9+0.5\\
&=45.78+0.5\\
&=46.28\end{aligned}}

\MethodBlock{方法三：按相同计数单位计算}{上方竖式及题目要求的验算仍保留；本法用于核对。先看成若干个百分之一，按整数加减，最后还原成小数。}{\begin{aligned}[t]&\phantom{={}}36.78+9.5\\
&=(3678+950)\div 100\\
&=4628\div 100\\
&=46.28\end{aligned}}
\par\noindent

\Answer{
\QHeading{计算过程：2. 列竖式计算}

小数点对齐，缺位补零；从右往左逐列计算。下面完整列出进位、退位及各数位的计算过程。

\par\medskip\noindent
\textbf{第1小题：\(50-14.14=35.86\)}\par\smallskip
\noindent
\begin{tikzpicture}[x=1em,y=1em,baseline=(current bounding box.center),draw=answerred,text=answerred]
\node[anchor=base,inner sep=0pt] at (0.000,0) {$5$};
\node[anchor=base,inner sep=0pt] at (0.800,0) {$0$};
\node[anchor=base,inner sep=0pt] at (1.600,0) {$.$};
\node[anchor=base,inner sep=0pt] at (1.900,0) {$0$};
\node[anchor=base,inner sep=0pt] at (2.700,0) {$0$};
\node[anchor=base,inner sep=0pt] at (0.000,-1.1) {$1$};
\node[anchor=base,inner sep=0pt] at (0.800,-1.1) {$4$};
\node[anchor=base,inner sep=0pt] at (1.600,-1.1) {$.$};
\node[anchor=base,inner sep=0pt] at (1.900,-1.1) {$1$};
\node[anchor=base,inner sep=0pt] at (2.700,-1.1) {$4$};
\node[anchor=base,inner sep=0pt] at (0.000,-2.3) {$3$};
\node[anchor=base,inner sep=0pt] at (0.800,-2.3) {$5$};
\node[anchor=base,inner sep=0pt] at (1.600,-2.3) {$.$};
\node[anchor=base,inner sep=0pt] at (1.900,-2.3) {$8$};
\node[anchor=base,inner sep=0pt] at (2.700,-2.3) {$6$};
\node[anchor=base,inner sep=0pt] at (-.7,-1.1) {$-$};
\draw[line width=.5pt] (-1,-1.52)--(3.300,-1.52);
\draw[line width=.35pt] (-0.160,.03)--(0.160,.43);
\node[anchor=base,inner sep=0pt,font=\normalsize] at (0.000,.72) {$4$};
\draw[line width=.35pt] (0.640,.03)--(0.960,.43);
\node[anchor=base,inner sep=0pt,font=\normalsize] at (0.800,.72) {$9$};
\draw[line width=.35pt] (1.740,.03)--(2.060,.43);
\node[anchor=base,inner sep=0pt,font=\normalsize] at (1.900,.72) {$9$};
\draw[line width=.35pt] (2.540,.03)--(2.860,.43);
\node[anchor=base,inner sep=0pt,font=\normalsize] at (2.700,.72) {$10$};
\end{tikzpicture}
\par\smallskip{\normalsize 划线数字是退位前的数，上方小字是退位后参与计算的数。\par}
\vspace{0pt}
\begin{enumerate}[leftmargin=1.7em,itemsep=5pt,topsep=0pt]
\item 十位退位：\(\begin{aligned}[t]&\phantom{={}}5-1\\&=4\end{aligned}\)\quad 借出1给后一位。
\item 个位传递退位：\(\begin{aligned}[t]&\phantom{={}}0+10-1\\&=10-1\\&=9\end{aligned}\)\quad 收到10，再借出1给后一位。
\item 十分位传递退位：\(\begin{aligned}[t]&\phantom{={}}0+10-1\\&=10-1\\&=9\end{aligned}\)\quad 收到10，再借出1给后一位。
\item 百分位接收退位：\(\begin{aligned}[t]&\phantom{={}}0+10\\&=10\end{aligned}\)\quad 前一位的1变成本位的10。
\item 百分位：\(\begin{aligned}[t]&\phantom{={}}10-4\\&=6\end{aligned}\)\quad 本位写6。
\item 十分位：\(\begin{aligned}[t]&\phantom{={}}9-1\\&=8\end{aligned}\)\quad 本位写8。
\item 个位：\(\begin{aligned}[t]&\phantom{={}}9-4\\&=5\end{aligned}\)\quad 本位写5。
\item 十位：\(\begin{aligned}[t]&\phantom{={}}4-1\\&=3\end{aligned}\)\quad 本位写3。
\end{enumerate}
\par\smallskip
因此，\(\begin{aligned}[t]&\phantom{={}}50-14.14\\&=35.86\end{aligned}\)。
\par\medskip

\par\medskip\noindent
\textbf{第2小题：\(53.69+14.86=68.55\)}\par\smallskip
\noindent
\begin{tikzpicture}[x=1em,y=1em,baseline=(current bounding box.center),draw=answerred,text=answerred]
\node[anchor=base,inner sep=0pt] at (0.000,0) {$5$};
\node[anchor=base,inner sep=0pt] at (0.800,0) {$3$};
\node[anchor=base,inner sep=0pt] at (1.600,0) {$.$};
\node[anchor=base,inner sep=0pt] at (1.900,0) {$6$};
\node[anchor=base,inner sep=0pt] at (2.700,0) {$9$};
\node[anchor=base,inner sep=0pt] at (0.000,-1.1) {$1$};
\node[anchor=base,inner sep=0pt] at (0.800,-1.1) {$4$};
\node[anchor=base,inner sep=0pt] at (1.600,-1.1) {$.$};
\node[anchor=base,inner sep=0pt] at (1.900,-1.1) {$8$};
\node[anchor=base,inner sep=0pt] at (2.700,-1.1) {$6$};
\node[anchor=base,inner sep=0pt] at (0.000,-2.3) {$6$};
\node[anchor=base,inner sep=0pt] at (0.800,-2.3) {$8$};
\node[anchor=base,inner sep=0pt] at (1.600,-2.3) {$.$};
\node[anchor=base,inner sep=0pt] at (1.900,-2.3) {$5$};
\node[anchor=base,inner sep=0pt] at (2.700,-2.3) {$5$};
\node[anchor=base,inner sep=0pt] at (-.7,-1.1) {$+$};
\draw[line width=.5pt] (-1,-1.52)--(3.300,-1.52);
\node[anchor=base,inner sep=0pt,font=\normalsize] at (1.900,.72) {$1$};
\node[anchor=base,inner sep=0pt,font=\normalsize] at (0.800,.72) {$1$};
\end{tikzpicture}
\par\smallskip{\normalsize 上方小字表示从后一位进来的数。\par}
\vspace{0pt}
\begin{enumerate}[leftmargin=1.7em,itemsep=5pt,topsep=0pt]
\item 百分位：\(\begin{aligned}[t]&\phantom{={}}9+6\\&=15\end{aligned}\)\quad 本位写5，向前一位进1。
\item 十分位：\(\begin{aligned}[t]&\phantom{={}}6+8+1\\&=14+1\\&=15\end{aligned}\)\quad 本位写5，向前一位进1。
\item 个位：\(\begin{aligned}[t]&\phantom{={}}3+4+1\\&=7+1\\&=8\end{aligned}\)\quad 本位写8。
\item 十位：\(\begin{aligned}[t]&\phantom{={}}5+1\\&=6\end{aligned}\)\quad 本位写6。
\end{enumerate}
\par\smallskip
因此，\(\begin{aligned}[t]&\phantom{={}}53.69+14.86\\&=68.55\end{aligned}\)。
\par\medskip

\par\medskip\noindent
\textbf{第3小题：\(53.62+12.59=66.21\)}\par\smallskip
\noindent
\begin{tikzpicture}[x=1em,y=1em,baseline=(current bounding box.center),draw=answerred,text=answerred]
\node[anchor=base,inner sep=0pt] at (0.000,0) {$5$};
\node[anchor=base,inner sep=0pt] at (0.800,0) {$3$};
\node[anchor=base,inner sep=0pt] at (1.600,0) {$.$};
\node[anchor=base,inner sep=0pt] at (1.900,0) {$6$};
\node[anchor=base,inner sep=0pt] at (2.700,0) {$2$};
\node[anchor=base,inner sep=0pt] at (0.000,-1.1) {$1$};
\node[anchor=base,inner sep=0pt] at (0.800,-1.1) {$2$};
\node[anchor=base,inner sep=0pt] at (1.600,-1.1) {$.$};
\node[anchor=base,inner sep=0pt] at (1.900,-1.1) {$5$};
\node[anchor=base,inner sep=0pt] at (2.700,-1.1) {$9$};
\node[anchor=base,inner sep=0pt] at (0.000,-2.3) {$6$};
\node[anchor=base,inner sep=0pt] at (0.800,-2.3) {$6$};
\node[anchor=base,inner sep=0pt] at (1.600,-2.3) {$.$};
\node[anchor=base,inner sep=0pt] at (1.900,-2.3) {$2$};
\node[anchor=base,inner sep=0pt] at (2.700,-2.3) {$1$};
\node[anchor=base,inner sep=0pt] at (-.7,-1.1) {$+$};
\draw[line width=.5pt] (-1,-1.52)--(3.300,-1.52);
\node[anchor=base,inner sep=0pt,font=\normalsize] at (1.900,.72) {$1$};
\node[anchor=base,inner sep=0pt,font=\normalsize] at (0.800,.72) {$1$};
\end{tikzpicture}
\par\smallskip{\normalsize 上方小字表示从后一位进来的数。\par}
\vspace{0pt}
\begin{enumerate}[leftmargin=1.7em,itemsep=5pt,topsep=0pt]
\item 百分位：\(\begin{aligned}[t]&\phantom{={}}2+9\\&=11\end{aligned}\)\quad 本位写1，向前一位进1。
\item 十分位：\(\begin{aligned}[t]&\phantom{={}}6+5+1\\&=11+1\\&=12\end{aligned}\)\quad 本位写2，向前一位进1。
\item 个位：\(\begin{aligned}[t]&\phantom{={}}3+2+1\\&=5+1\\&=6\end{aligned}\)\quad 本位写6。
\item 十位：\(\begin{aligned}[t]&\phantom{={}}5+1\\&=6\end{aligned}\)\quad 本位写6。
\end{enumerate}
\par\smallskip
因此，\(\begin{aligned}[t]&\phantom{={}}53.62+12.59\\&=66.21\end{aligned}\)。
\par\medskip

\par\medskip\noindent
\textbf{第4小题：\(12.98-3.8=9.18\)}\par\smallskip
\noindent
\begin{tikzpicture}[x=1em,y=1em,baseline=(current bounding box.center),draw=answerred,text=answerred]
\node[anchor=base,inner sep=0pt] at (0.000,0) {$1$};
\node[anchor=base,inner sep=0pt] at (0.800,0) {$2$};
\node[anchor=base,inner sep=0pt] at (1.600,0) {$.$};
\node[anchor=base,inner sep=0pt] at (1.900,0) {$9$};
\node[anchor=base,inner sep=0pt] at (2.700,0) {$8$};
\node[anchor=base,inner sep=0pt] at (0.800,-1.1) {$3$};
\node[anchor=base,inner sep=0pt] at (1.600,-1.1) {$.$};
\node[anchor=base,inner sep=0pt] at (1.900,-1.1) {$8$};
\node[anchor=base,inner sep=0pt] at (2.700,-1.1) {$0$};
\node[anchor=base,inner sep=0pt] at (0.800,-2.3) {$9$};
\node[anchor=base,inner sep=0pt] at (1.600,-2.3) {$.$};
\node[anchor=base,inner sep=0pt] at (1.900,-2.3) {$1$};
\node[anchor=base,inner sep=0pt] at (2.700,-2.3) {$8$};
\node[anchor=base,inner sep=0pt] at (-.7,-1.1) {$-$};
\draw[line width=.5pt] (-1,-1.52)--(3.300,-1.52);
\draw[line width=.35pt] (-0.160,.03)--(0.160,.43);
\node[anchor=base,inner sep=0pt,font=\normalsize] at (0.000,.72) {$0$};
\draw[line width=.35pt] (0.640,.03)--(0.960,.43);
\node[anchor=base,inner sep=0pt,font=\normalsize] at (0.800,.72) {$12$};
\end{tikzpicture}
\par\smallskip{\normalsize 划线数字是退位前的数，上方小字是退位后参与计算的数。\par}
\vspace{0pt}
\begin{enumerate}[leftmargin=1.7em,itemsep=5pt,topsep=0pt]
\item 百分位：\(\begin{aligned}[t]&\phantom{={}}8-0\\&=8\end{aligned}\)\quad 本位写8。
\item 十分位：\(\begin{aligned}[t]&\phantom{={}}9-8\\&=1\end{aligned}\)\quad 本位写1。
\item 十位退位：\(\begin{aligned}[t]&\phantom{={}}1-1\\&=0\end{aligned}\)\quad 借出1给后一位。
\item 个位接收退位：\(\begin{aligned}[t]&\phantom{={}}2+10\\&=12\end{aligned}\)\quad 前一位的1变成本位的10。
\item 个位：\(\begin{aligned}[t]&\phantom{={}}12-3\\&=9\end{aligned}\)\quad 本位写9。
\item 十位：\(\begin{aligned}[t]&\phantom{={}}0-0\\&=0\end{aligned}\)\quad 本位写0。
\end{enumerate}
\par\smallskip
因此，\(\begin{aligned}[t]&\phantom{={}}12.98-3.8\\&=9.18\end{aligned}\)。
\par\medskip

\par\medskip\noindent
\textbf{第5小题：\(12.96-6.2=6.76\)}\par\smallskip
\noindent
\begin{tikzpicture}[x=1em,y=1em,baseline=(current bounding box.center),draw=answerred,text=answerred]
\node[anchor=base,inner sep=0pt] at (0.000,0) {$1$};
\node[anchor=base,inner sep=0pt] at (0.800,0) {$2$};
\node[anchor=base,inner sep=0pt] at (1.600,0) {$.$};
\node[anchor=base,inner sep=0pt] at (1.900,0) {$9$};
\node[anchor=base,inner sep=0pt] at (2.700,0) {$6$};
\node[anchor=base,inner sep=0pt] at (0.800,-1.1) {$6$};
\node[anchor=base,inner sep=0pt] at (1.600,-1.1) {$.$};
\node[anchor=base,inner sep=0pt] at (1.900,-1.1) {$2$};
\node[anchor=base,inner sep=0pt] at (2.700,-1.1) {$0$};
\node[anchor=base,inner sep=0pt] at (0.800,-2.3) {$6$};
\node[anchor=base,inner sep=0pt] at (1.600,-2.3) {$.$};
\node[anchor=base,inner sep=0pt] at (1.900,-2.3) {$7$};
\node[anchor=base,inner sep=0pt] at (2.700,-2.3) {$6$};
\node[anchor=base,inner sep=0pt] at (-.7,-1.1) {$-$};
\draw[line width=.5pt] (-1,-1.52)--(3.300,-1.52);
\draw[line width=.35pt] (-0.160,.03)--(0.160,.43);
\node[anchor=base,inner sep=0pt,font=\normalsize] at (0.000,.72) {$0$};
\draw[line width=.35pt] (0.640,.03)--(0.960,.43);
\node[anchor=base,inner sep=0pt,font=\normalsize] at (0.800,.72) {$12$};
\end{tikzpicture}
\par\smallskip{\normalsize 划线数字是退位前的数，上方小字是退位后参与计算的数。\par}
\vspace{0pt}
\begin{enumerate}[leftmargin=1.7em,itemsep=5pt,topsep=0pt]
\item 百分位：\(\begin{aligned}[t]&\phantom{={}}6-0\\&=6\end{aligned}\)\quad 本位写6。
\item 十分位：\(\begin{aligned}[t]&\phantom{={}}9-2\\&=7\end{aligned}\)\quad 本位写7。
\item 十位退位：\(\begin{aligned}[t]&\phantom{={}}1-1\\&=0\end{aligned}\)\quad 借出1给后一位。
\item 个位接收退位：\(\begin{aligned}[t]&\phantom{={}}2+10\\&=12\end{aligned}\)\quad 前一位的1变成本位的10。
\item 个位：\(\begin{aligned}[t]&\phantom{={}}12-6\\&=6\end{aligned}\)\quad 本位写6。
\item 十位：\(\begin{aligned}[t]&\phantom{={}}0-0\\&=0\end{aligned}\)\quad 本位写0。
\end{enumerate}
\par\smallskip
因此，\(\begin{aligned}[t]&\phantom{={}}12.96-6.2\\&=6.76\end{aligned}\)。
\par\medskip

\par\medskip\noindent
\textbf{第6小题：\(4.05+9.09=13.14\)}\par\smallskip
\noindent
\begin{tikzpicture}[x=1em,y=1em,baseline=(current bounding box.center),draw=answerred,text=answerred]
\node[anchor=base,inner sep=0pt] at (0.800,0) {$4$};
\node[anchor=base,inner sep=0pt] at (1.600,0) {$.$};
\node[anchor=base,inner sep=0pt] at (1.900,0) {$0$};
\node[anchor=base,inner sep=0pt] at (2.700,0) {$5$};
\node[anchor=base,inner sep=0pt] at (0.800,-1.1) {$9$};
\node[anchor=base,inner sep=0pt] at (1.600,-1.1) {$.$};
\node[anchor=base,inner sep=0pt] at (1.900,-1.1) {$0$};
\node[anchor=base,inner sep=0pt] at (2.700,-1.1) {$9$};
\node[anchor=base,inner sep=0pt] at (0.000,-2.3) {$1$};
\node[anchor=base,inner sep=0pt] at (0.800,-2.3) {$3$};
\node[anchor=base,inner sep=0pt] at (1.600,-2.3) {$.$};
\node[anchor=base,inner sep=0pt] at (1.900,-2.3) {$1$};
\node[anchor=base,inner sep=0pt] at (2.700,-2.3) {$4$};
\node[anchor=base,inner sep=0pt] at (-.7,-1.1) {$+$};
\draw[line width=.5pt] (-1,-1.52)--(3.300,-1.52);
\node[anchor=base,inner sep=0pt,font=\normalsize] at (1.900,.72) {$1$};
\node[anchor=base,inner sep=0pt,font=\normalsize] at (0.000,.72) {$1$};
\end{tikzpicture}
\par\smallskip{\normalsize 上方小字表示从后一位进来的数。\par}
\vspace{0pt}
\begin{enumerate}[leftmargin=1.7em,itemsep=5pt,topsep=0pt]
\item 百分位：\(\begin{aligned}[t]&\phantom{={}}5+9\\&=14\end{aligned}\)\quad 本位写4，向前一位进1。
\item 十分位：\(\begin{aligned}[t]&\phantom{={}}0+0+1\\&=0+1\\&=1\end{aligned}\)\quad 本位写1。
\item 个位：\(\begin{aligned}[t]&\phantom{={}}4+9\\&=13\end{aligned}\)\quad 本位写3，向前一位进1。
\end{enumerate}
最高位再写进上来的1。\par
\par\smallskip
因此，\(\begin{aligned}[t]&\phantom{={}}4.05+9.09\\&=13.14\end{aligned}\)。
\par\medskip

\par\medskip\noindent
\textbf{第7小题：\(2.56+4.44=7\)}\par\smallskip
\noindent
\begin{tikzpicture}[x=1em,y=1em,baseline=(current bounding box.center),draw=answerred,text=answerred]
\node[anchor=base,inner sep=0pt] at (0.000,0) {$2$};
\node[anchor=base,inner sep=0pt] at (0.800,0) {$.$};
\node[anchor=base,inner sep=0pt] at (1.100,0) {$5$};
\node[anchor=base,inner sep=0pt] at (1.900,0) {$6$};
\node[anchor=base,inner sep=0pt] at (0.000,-1.1) {$4$};
\node[anchor=base,inner sep=0pt] at (0.800,-1.1) {$.$};
\node[anchor=base,inner sep=0pt] at (1.100,-1.1) {$4$};
\node[anchor=base,inner sep=0pt] at (1.900,-1.1) {$4$};
\node[anchor=base,inner sep=0pt] at (0.000,-2.3) {$7$};
\node[anchor=base,inner sep=0pt] at (0.800,-2.3) {$.$};
\node[anchor=base,inner sep=0pt] at (1.100,-2.3) {$0$};
\node[anchor=base,inner sep=0pt] at (1.900,-2.3) {$0$};
\node[anchor=base,inner sep=0pt] at (-.7,-1.1) {$+$};
\draw[line width=.5pt] (-1,-1.52)--(2.500,-1.52);
\node[anchor=base,inner sep=0pt,font=\normalsize] at (1.100,.72) {$1$};
\node[anchor=base,inner sep=0pt,font=\normalsize] at (0.000,.72) {$1$};
\end{tikzpicture}
\par\smallskip{\normalsize 上方小字表示从后一位进来的数。\par}
\vspace{0pt}
\begin{enumerate}[leftmargin=1.7em,itemsep=5pt,topsep=0pt]
\item 百分位：\(\begin{aligned}[t]&\phantom{={}}6+4\\&=10\end{aligned}\)\quad 本位写0，向前一位进1。
\item 十分位：\(\begin{aligned}[t]&\phantom{={}}5+4+1\\&=9+1\\&=10\end{aligned}\)\quad 本位写0，向前一位进1。
\item 个位：\(\begin{aligned}[t]&\phantom{={}}2+4+1\\&=6+1\\&=7\end{aligned}\)\quad 本位写7。
\end{enumerate}
\par\smallskip
因此，\(\begin{aligned}[t]&\phantom{={}}2.56+4.44\\&=7.00\end{aligned}\)。
\par\medskip

\par\medskip\noindent
\textbf{第7小题验算：\(7.00-4.44\)}\par\smallskip
\noindent
\begin{tikzpicture}[x=1em,y=1em,baseline=(current bounding box.center),draw=answerred,text=answerred]
\node[anchor=base,inner sep=0pt] at (0.000,0) {$7$};
\node[anchor=base,inner sep=0pt] at (0.800,0) {$.$};
\node[anchor=base,inner sep=0pt] at (1.100,0) {$0$};
\node[anchor=base,inner sep=0pt] at (1.900,0) {$0$};
\node[anchor=base,inner sep=0pt] at (0.000,-1.1) {$4$};
\node[anchor=base,inner sep=0pt] at (0.800,-1.1) {$.$};
\node[anchor=base,inner sep=0pt] at (1.100,-1.1) {$4$};
\node[anchor=base,inner sep=0pt] at (1.900,-1.1) {$4$};
\node[anchor=base,inner sep=0pt] at (0.000,-2.3) {$2$};
\node[anchor=base,inner sep=0pt] at (0.800,-2.3) {$.$};
\node[anchor=base,inner sep=0pt] at (1.100,-2.3) {$5$};
\node[anchor=base,inner sep=0pt] at (1.900,-2.3) {$6$};
\node[anchor=base,inner sep=0pt] at (-.7,-1.1) {$-$};
\draw[line width=.5pt] (-1,-1.52)--(2.500,-1.52);
\draw[line width=.35pt] (-0.160,.03)--(0.160,.43);
\node[anchor=base,inner sep=0pt,font=\normalsize] at (0.000,.72) {$6$};
\draw[line width=.35pt] (0.940,.03)--(1.260,.43);
\node[anchor=base,inner sep=0pt,font=\normalsize] at (1.100,.72) {$9$};
\draw[line width=.35pt] (1.740,.03)--(2.060,.43);
\node[anchor=base,inner sep=0pt,font=\normalsize] at (1.900,.72) {$10$};
\end{tikzpicture}
\par\smallskip{\normalsize 划线数字是退位前的数，上方小字是退位后参与计算的数。\par}
\vspace{0pt}
\begin{enumerate}[leftmargin=1.7em,itemsep=5pt,topsep=0pt]
\item 个位退位：\(\begin{aligned}[t]&\phantom{={}}7-1\\&=6\end{aligned}\)\quad 借出1给后一位。
\item 十分位传递退位：\(\begin{aligned}[t]&\phantom{={}}0+10-1\\&=10-1\\&=9\end{aligned}\)\quad 收到10，再借出1给后一位。
\item 百分位接收退位：\(\begin{aligned}[t]&\phantom{={}}0+10\\&=10\end{aligned}\)\quad 前一位的1变成本位的10。
\item 百分位：\(\begin{aligned}[t]&\phantom{={}}10-4\\&=6\end{aligned}\)\quad 本位写6。
\item 十分位：\(\begin{aligned}[t]&\phantom{={}}9-4\\&=5\end{aligned}\)\quad 本位写5。
\item 个位：\(\begin{aligned}[t]&\phantom{={}}6-4\\&=2\end{aligned}\)\quad 本位写2。
\end{enumerate}
\par\smallskip
因此，\(\begin{aligned}[t]&\phantom{={}}7.00-4.44\\&=2.56\end{aligned}\)。
\par\medskip

\par\medskip\noindent
\textbf{第8小题：\(28.3-15.06=13.24\)}\par\smallskip
\noindent
\begin{tikzpicture}[x=1em,y=1em,baseline=(current bounding box.center),draw=answerred,text=answerred]
\node[anchor=base,inner sep=0pt] at (0.000,0) {$2$};
\node[anchor=base,inner sep=0pt] at (0.800,0) {$8$};
\node[anchor=base,inner sep=0pt] at (1.600,0) {$.$};
\node[anchor=base,inner sep=0pt] at (1.900,0) {$3$};
\node[anchor=base,inner sep=0pt] at (2.700,0) {$0$};
\node[anchor=base,inner sep=0pt] at (0.000,-1.1) {$1$};
\node[anchor=base,inner sep=0pt] at (0.800,-1.1) {$5$};
\node[anchor=base,inner sep=0pt] at (1.600,-1.1) {$.$};
\node[anchor=base,inner sep=0pt] at (1.900,-1.1) {$0$};
\node[anchor=base,inner sep=0pt] at (2.700,-1.1) {$6$};
\node[anchor=base,inner sep=0pt] at (0.000,-2.3) {$1$};
\node[anchor=base,inner sep=0pt] at (0.800,-2.3) {$3$};
\node[anchor=base,inner sep=0pt] at (1.600,-2.3) {$.$};
\node[anchor=base,inner sep=0pt] at (1.900,-2.3) {$2$};
\node[anchor=base,inner sep=0pt] at (2.700,-2.3) {$4$};
\node[anchor=base,inner sep=0pt] at (-.7,-1.1) {$-$};
\draw[line width=.5pt] (-1,-1.52)--(3.300,-1.52);
\draw[line width=.35pt] (1.740,.03)--(2.060,.43);
\node[anchor=base,inner sep=0pt,font=\normalsize] at (1.900,.72) {$2$};
\draw[line width=.35pt] (2.540,.03)--(2.860,.43);
\node[anchor=base,inner sep=0pt,font=\normalsize] at (2.700,.72) {$10$};
\end{tikzpicture}
\par\smallskip{\normalsize 划线数字是退位前的数，上方小字是退位后参与计算的数。\par}
\vspace{0pt}
\begin{enumerate}[leftmargin=1.7em,itemsep=5pt,topsep=0pt]
\item 十分位退位：\(\begin{aligned}[t]&\phantom{={}}3-1\\&=2\end{aligned}\)\quad 借出1给后一位。
\item 百分位接收退位：\(\begin{aligned}[t]&\phantom{={}}0+10\\&=10\end{aligned}\)\quad 前一位的1变成本位的10。
\item 百分位：\(\begin{aligned}[t]&\phantom{={}}10-6\\&=4\end{aligned}\)\quad 本位写4。
\item 十分位：\(\begin{aligned}[t]&\phantom{={}}2-0\\&=2\end{aligned}\)\quad 本位写2。
\item 个位：\(\begin{aligned}[t]&\phantom{={}}8-5\\&=3\end{aligned}\)\quad 本位写3。
\item 十位：\(\begin{aligned}[t]&\phantom{={}}2-1\\&=1\end{aligned}\)\quad 本位写1。
\end{enumerate}
\par\smallskip
因此，\(\begin{aligned}[t]&\phantom{={}}28.3-15.06\\&=13.24\end{aligned}\)。
\par\medskip

\par\medskip\noindent
\textbf{第8小题验算：\(13.24+15.06\)}\par\smallskip
\noindent
\begin{tikzpicture}[x=1em,y=1em,baseline=(current bounding box.center),draw=answerred,text=answerred]
\node[anchor=base,inner sep=0pt] at (0.000,0) {$1$};
\node[anchor=base,inner sep=0pt] at (0.800,0) {$3$};
\node[anchor=base,inner sep=0pt] at (1.600,0) {$.$};
\node[anchor=base,inner sep=0pt] at (1.900,0) {$2$};
\node[anchor=base,inner sep=0pt] at (2.700,0) {$4$};
\node[anchor=base,inner sep=0pt] at (0.000,-1.1) {$1$};
\node[anchor=base,inner sep=0pt] at (0.800,-1.1) {$5$};
\node[anchor=base,inner sep=0pt] at (1.600,-1.1) {$.$};
\node[anchor=base,inner sep=0pt] at (1.900,-1.1) {$0$};
\node[anchor=base,inner sep=0pt] at (2.700,-1.1) {$6$};
\node[anchor=base,inner sep=0pt] at (0.000,-2.3) {$2$};
\node[anchor=base,inner sep=0pt] at (0.800,-2.3) {$8$};
\node[anchor=base,inner sep=0pt] at (1.600,-2.3) {$.$};
\node[anchor=base,inner sep=0pt] at (1.900,-2.3) {$3$};
\node[anchor=base,inner sep=0pt] at (2.700,-2.3) {$0$};
\node[anchor=base,inner sep=0pt] at (-.7,-1.1) {$+$};
\draw[line width=.5pt] (-1,-1.52)--(3.300,-1.52);
\node[anchor=base,inner sep=0pt,font=\normalsize] at (1.900,.72) {$1$};
\end{tikzpicture}
\par\smallskip{\normalsize 上方小字表示从后一位进来的数。\par}
\vspace{0pt}
\begin{enumerate}[leftmargin=1.7em,itemsep=5pt,topsep=0pt]
\item 百分位：\(\begin{aligned}[t]&\phantom{={}}4+6\\&=10\end{aligned}\)\quad 本位写0，向前一位进1。
\item 十分位：\(\begin{aligned}[t]&\phantom{={}}2+0+1\\&=2+1\\&=3\end{aligned}\)\quad 本位写3。
\item 个位：\(\begin{aligned}[t]&\phantom{={}}3+5\\&=8\end{aligned}\)\quad 本位写8。
\item 十位：\(\begin{aligned}[t]&\phantom{={}}1+1\\&=2\end{aligned}\)\quad 本位写2。
\end{enumerate}
\par\smallskip
因此，\(\begin{aligned}[t]&\phantom{={}}13.24+15.06\\&=28.30\end{aligned}\)。
\par\medskip

\par\medskip\noindent
\textbf{第9小题：\(36.78+9.5=46.28\)}\par\smallskip
\noindent
\begin{tikzpicture}[x=1em,y=1em,baseline=(current bounding box.center),draw=answerred,text=answerred]
\node[anchor=base,inner sep=0pt] at (0.000,0) {$3$};
\node[anchor=base,inner sep=0pt] at (0.800,0) {$6$};
\node[anchor=base,inner sep=0pt] at (1.600,0) {$.$};
\node[anchor=base,inner sep=0pt] at (1.900,0) {$7$};
\node[anchor=base,inner sep=0pt] at (2.700,0) {$8$};
\node[anchor=base,inner sep=0pt] at (0.800,-1.1) {$9$};
\node[anchor=base,inner sep=0pt] at (1.600,-1.1) {$.$};
\node[anchor=base,inner sep=0pt] at (1.900,-1.1) {$5$};
\node[anchor=base,inner sep=0pt] at (2.700,-1.1) {$0$};
\node[anchor=base,inner sep=0pt] at (0.000,-2.3) {$4$};
\node[anchor=base,inner sep=0pt] at (0.800,-2.3) {$6$};
\node[anchor=base,inner sep=0pt] at (1.600,-2.3) {$.$};
\node[anchor=base,inner sep=0pt] at (1.900,-2.3) {$2$};
\node[anchor=base,inner sep=0pt] at (2.700,-2.3) {$8$};
\node[anchor=base,inner sep=0pt] at (-.7,-1.1) {$+$};
\draw[line width=.5pt] (-1,-1.52)--(3.300,-1.52);
\node[anchor=base,inner sep=0pt,font=\normalsize] at (0.800,.72) {$1$};
\node[anchor=base,inner sep=0pt,font=\normalsize] at (0.000,.72) {$1$};
\end{tikzpicture}
\par\smallskip{\normalsize 上方小字表示从后一位进来的数。\par}
\vspace{0pt}
\begin{enumerate}[leftmargin=1.7em,itemsep=5pt,topsep=0pt]
\item 百分位：\(\begin{aligned}[t]&\phantom{={}}8+0\\&=8\end{aligned}\)\quad 本位写8。
\item 十分位：\(\begin{aligned}[t]&\phantom{={}}7+5\\&=12\end{aligned}\)\quad 本位写2，向前一位进1。
\item 个位：\(\begin{aligned}[t]&\phantom{={}}6+9+1\\&=15+1\\&=16\end{aligned}\)\quad 本位写6，向前一位进1。
\item 十位：\(\begin{aligned}[t]&\phantom{={}}3+0+1\\&=3+1\\&=4\end{aligned}\)\quad 本位写4。
\end{enumerate}
\par\smallskip
因此，\(\begin{aligned}[t]&\phantom{={}}36.78+9.5\\&=46.28\end{aligned}\)。
\par\medskip

\par\medskip\noindent
\textbf{第9小题验算：\(46.28-9.5\)}\par\smallskip
\noindent
\begin{tikzpicture}[x=1em,y=1em,baseline=(current bounding box.center),draw=answerred,text=answerred]
\node[anchor=base,inner sep=0pt] at (0.000,0) {$4$};
\node[anchor=base,inner sep=0pt] at (0.800,0) {$6$};
\node[anchor=base,inner sep=0pt] at (1.600,0) {$.$};
\node[anchor=base,inner sep=0pt] at (1.900,0) {$2$};
\node[anchor=base,inner sep=0pt] at (2.700,0) {$8$};
\node[anchor=base,inner sep=0pt] at (0.800,-1.1) {$9$};
\node[anchor=base,inner sep=0pt] at (1.600,-1.1) {$.$};
\node[anchor=base,inner sep=0pt] at (1.900,-1.1) {$5$};
\node[anchor=base,inner sep=0pt] at (2.700,-1.1) {$0$};
\node[anchor=base,inner sep=0pt] at (0.000,-2.3) {$3$};
\node[anchor=base,inner sep=0pt] at (0.800,-2.3) {$6$};
\node[anchor=base,inner sep=0pt] at (1.600,-2.3) {$.$};
\node[anchor=base,inner sep=0pt] at (1.900,-2.3) {$7$};
\node[anchor=base,inner sep=0pt] at (2.700,-2.3) {$8$};
\node[anchor=base,inner sep=0pt] at (-.7,-1.1) {$-$};
\draw[line width=.5pt] (-1,-1.52)--(3.300,-1.52);
\draw[line width=.35pt] (-0.160,.03)--(0.160,.43);
\node[anchor=base,inner sep=0pt,font=\normalsize] at (0.000,.72) {$3$};
\draw[line width=.35pt] (0.640,.03)--(0.960,.43);
\node[anchor=base,inner sep=0pt,font=\normalsize] at (0.800,.72) {$15$};
\draw[line width=.35pt] (1.740,.03)--(2.060,.43);
\node[anchor=base,inner sep=0pt,font=\normalsize] at (1.900,.72) {$12$};
\end{tikzpicture}
\par\smallskip{\normalsize 划线数字是退位前的数，上方小字是退位后参与计算的数。\par}
\vspace{0pt}
\begin{enumerate}[leftmargin=1.7em,itemsep=5pt,topsep=0pt]
\item 百分位：\(\begin{aligned}[t]&\phantom{={}}8-0\\&=8\end{aligned}\)\quad 本位写8。
\item 个位退位：\(\begin{aligned}[t]&\phantom{={}}6-1\\&=5\end{aligned}\)\quad 借出1给后一位。
\item 十分位接收退位：\(\begin{aligned}[t]&\phantom{={}}2+10\\&=12\end{aligned}\)\quad 前一位的1变成本位的10。
\item 十分位：\(\begin{aligned}[t]&\phantom{={}}12-5\\&=7\end{aligned}\)\quad 本位写7。
\item 十位退位：\(\begin{aligned}[t]&\phantom{={}}4-1\\&=3\end{aligned}\)\quad 借出1给后一位。
\item 个位接收退位：\(\begin{aligned}[t]&\phantom{={}}5+10\\&=15\end{aligned}\)\quad 前一位的1变成本位的10。
\item 个位：\(\begin{aligned}[t]&\phantom{={}}15-9\\&=6\end{aligned}\)\quad 本位写6。
\item 十位：\(\begin{aligned}[t]&\phantom{={}}3-0\\&=3\end{aligned}\)\quad 本位写3。
\end{enumerate}
\par\smallskip
因此，\(\begin{aligned}[t]&\phantom{={}}46.28-9.5\\&=36.78\end{aligned}\)。
\par\medskip

}

\QHeading{第3题　解决问题。}

\Problem{① 在一次跳高比赛中，亮亮跳过了1.15米，华华比亮亮跳的高0.08米，明明比华华跳的低0.16米。明明跳过了多少米？}
\par


\MethodTitle{方法一}
\Answer{{\Formula{\begin{aligned}[t]&\phantom{={}}1.15+0.08-0.16\\&=1.23-0.16\\&=1.07\text{（米）}\end{aligned}}}\par 答：明明跳过了1.07米。}

\MethodTitle{方法二：先求明明比亮亮低多少}
\Answer{明明比华华低0.16米，而华华比亮亮高0.08米，所以明明比亮亮低0.08米。}
\Answer{\Formula{\begin{aligned}[t]&\phantom{={}}1.15-(0.16-0.08)\\
&=1.15-0.08\\
&=1.07\text{（米）}\end{aligned}}}
\Answer{答：明明跳过了1.07米。}
\Problem{② 李伯伯从果园里摘了一些苹果去售卖，上午卖了12.8千克，下午比上午少卖了0.15千克，最后还剩18.15千克没卖完。李伯伯摘了多少千克苹果拿去售卖？}

\MethodTitle{方法一：简便计算}
\Answer{{\Formula{\begin{aligned}[t]&\phantom{={}}12.8+(12.8-0.15)+18.15\\&=12.8+12.8-0.15+18.15\\&=(12.8+12.8)+(18.15-0.15)\\&=25.6+18\\&=43.6\text{（千克）}\end{aligned}}}\par 答：原有43.6千克。}

\MethodTitle{方法二：先按两次都卖12.8千克计算}
\Answer{把下午也看成12.8千克，多算0.15千克；再加剩下的苹果。}
\Answer{\Formula{\begin{aligned}[t]&\phantom{={}}12.8\times 2-0.15+18.15\\
&=25.6-0.15+18.15\\
&=25.45+18.15\\
&=43.6\text{（千克）}\end{aligned}}}
\Answer{答：李伯伯摘了43.6千克苹果。}

